% Generated. Do not edit.
% Projection scriptural-date-location/section, version 1: the Scriptural Date and
% Location page, generated whole from research/chronology.toml, the Scripture
% chronology corpus and the source library by tools/tpt document-types refresh.
% document: liturgy/roman-rite/postconciliar/roman-missal-third-edition-en-us-2011/propers/general-calendar/pc-r07-all-saints-abc
\runninghead{Scriptural Date and Location}
\section{Scriptural Date and Location}
\label{sec:date-location}

{\small Each passage the Mass reads or sings from Scripture, in canonical order. The Context column names what is dated --- the event the passage narrates, a setting it remembers, the traditional attribution of its author, or the composition of its book --- and where it belongs; the Date column gives the answer of the Scripture chronology corpus for that subject, with every alternative it records. Texts of the Mass that are not drawn from Scripture have no row. The evidence for each date follows the table.\par}
\medskip
% Generated. Do not edit.
% Deterministic projection of research/chronology.toml.
% schema: 5
% document: liturgy/roman-rite/postconciliar/roman-missal-third-edition-en-us-2011/propers/general-calendar/pc-r07-all-saints-abc
% generated-by: tools/tpt proper-chronology annotations
\newcommand{\chronologyannotationclaim}[6]{#6}
\newcommand{\chronologyannotationreach}[3]{}
\newcommand{\chronologyannotationgroup}[3]{#3}
\newcommand{\chronologyannotation}[1]{%
  \ifcsname triptychchronologyannotation@#1\endcsname
    \csname triptychchronologyannotation@#1\endcsname
  \else
    \PackageError{triptych}{No chronology annotation for #1}{}%
  \fi}
% appointment-dependency: guidance/liturgy/postconciliar-propers-registry.md sha256=376860a7c37a4c75bb6d880cd1b260183be1877a07e3cd01660a9e2854e2ea72
% appointment-dependency: src/gpt/liturgy/roman-rite/postconciliar/roman-missal-third-edition-en-us-2011/propers/general-calendar/pc-r07-all-saints-abc/instance/manifest.md sha256=9ae6055704d1fe2aa3949d5e91e816673c113bc775b3a39b1bc145c4290b6630
% appointment-dependency: src/gpt/liturgy/roman-rite/postconciliar/roman-missal-third-edition-en-us-2011/propers/general-calendar/pc-r07-all-saints-abc/propers/verified.md sha256=574617c10630e751cb43a297cfec38b5c622631c3dd5ca4b8092d3f70c5441b5
% appointment-dependency: src/gpt/liturgy/roman-rite/postconciliar/roman-missal-third-edition-en-us-2011/propers/general-calendar/pc-r07-all-saints-abc/research/chronology-inputs.toml sha256=f30b32a467cb7d56012eddb3291e509def6fa3b786ee844bead09d3b8452e975
% appointment-dependency: src/gpt/liturgy/roman-rite/postconciliar/roman-missal-third-edition-en-us-2011/propers/general-calendar/shared/formularies/pc-r07-all-saints/propers/verified.md sha256=3025d42355a3162576f155c12ed1d0eadca32ec592546ef8ca6bf7c7066e5630
% appointment-dependency: src/gpt/liturgy/roman-rite/postconciliar/roman-missal-third-edition-en-us-2011/propers/registry/README.md sha256=df0aae579b2ae7e02cca5597f382ee615974bc84fc46fea909a9c21d942ce152
% appointment-dependency: src/gpt/liturgy/roman-rite/postconciliar/roman-missal-third-edition-en-us-2011/propers/registry/formula-dispositions.md sha256=af014f05723fb53a99006ef740101dc8abf629ed8b0b53a60069e94713cb289e
% appointment-note: entrance-antiphon: Composed Missal text without an appointed biblical citation; theological echoes and compiler chant directions do not create appointments.
% appointment-note: collect: Composed Missal text without an appointed biblical citation; theological echoes and compiler chant directions do not create appointments.
% appointment-note: first-reading: Revelation 7:2-4,9-14; appointed; vulgate numbering; canonical verse queries, not appointed wording.
% appointment-note: responsorial-psalm: Psalm 24:1bc-2,3-4ab,5-6; appointed; hebrew numbering; whole-verse query envelope, not subverse chronology or appointed wording.
% appointment-note: responsorial-psalm: Cf. Psalm 24:6; adaptation; hebrew numbering; canonical verse queries, not appointed wording.
% appointment-note: second-reading: 1 John 3:1-3; appointed; vulgate numbering; canonical verse queries, not appointed wording.
% appointment-note: gospel-acclamation: Matthew 11:28; appointed; vulgate numbering; canonical verse queries, not appointed wording.
% appointment-note: gospel: Matthew 5:1-12a; appointed; vulgate numbering; whole-verse query envelope, not subverse chronology or appointed wording.
% appointment-note: prayer-over-offerings: Composed Missal text without an appointed biblical citation; theological echoes and compiler chant directions do not create appointments.
% appointment-note: proper-preface: Composed Missal text without an appointed biblical citation; theological echoes and compiler chant directions do not create appointments.
% appointment-note: communion-antiphon: Matthew 5:8-10; appointed; vulgate numbering; canonical verse queries, not appointed wording.
% appointment-note: prayer-after-communion: Composed Missal text without an appointed biblical citation; theological echoes and compiler chant directions do not create appointments.
% appointment-note: solemn-blessing: Composed Missal text without an appointed biblical citation; theological echoes and compiler chant directions do not create appointments.
% first-reading: First Reading
\expandafter\def\csname triptychchronologyannotation@first-reading\endcsname{%
  \chronologyannotationgroup{composition}{disputed}{\textbf{Composition} -- disputed: \chronologyannotationreach{Apoc.7.10}{Apoc}{true}\chronologyannotationreach{Apoc.7.11}{Apoc}{true}\chronologyannotationreach{Apoc.7.12}{Apoc}{true}\chronologyannotationreach{Apoc.7.13}{Apoc}{true}\chronologyannotationreach{Apoc.7.14}{Apoc}{true}\chronologyannotationreach{Apoc.7.2}{Apoc}{true}\chronologyannotationreach{Apoc.7.3}{Apoc}{true}\chronologyannotationreach{Apoc.7.4}{Apoc}{true}\chronologyannotationreach{Apoc.7.9}{Apoc}{true}\chronologyannotationclaim{composition.apocalypse}{composition}{catholic-traditional-v1}{disputed}{the reign of Domitian ... A.D. 81-96}{A.D. 81--96}; \chronologyannotationreach{Apoc.7.10}{Apoc}{true}\chronologyannotationreach{Apoc.7.11}{Apoc}{true}\chronologyannotationreach{Apoc.7.12}{Apoc}{true}\chronologyannotationreach{Apoc.7.13}{Apoc}{true}\chronologyannotationreach{Apoc.7.14}{Apoc}{true}\chronologyannotationreach{Apoc.7.2}{Apoc}{true}\chronologyannotationreach{Apoc.7.3}{Apoc}{true}\chronologyannotationreach{Apoc.7.4}{Apoc}{true}\chronologyannotationreach{Apoc.7.9}{Apoc}{true}\chronologyannotationclaim{composition.apocalypse}{composition}{catholic-traditional-v1}{disputed}{from the year 90 to 100 (approximately)}{c. A.D. 90--100}.}%
}
% responsorial-psalm: Responsorial Psalm
\expandafter\def\csname triptychchronologyannotation@responsorial-psalm\endcsname{%
  \chronologyannotationgroup{historical-setting}{research-pending}{\textbf{Historical setting} -- Occasion date unresolved.}%
  \space \chronologyannotationgroup{composition}{preferred}{\textbf{Composition}: \chronologyannotationreach{Ps.23.1}{Ps}{true}\chronologyannotationreach{Ps.23.2}{Ps}{true}\chronologyannotationreach{Ps.23.3}{Ps}{true}\chronologyannotationreach{Ps.23.4}{Ps}{true}\chronologyannotationreach{Ps.23.5}{Ps}{true}\chronologyannotationreach{Ps.23.6}{Ps}{true}\chronologyannotationclaim{critical.psalms.latest-composition-boundary}{composition}{catholic-critical-v1}{preferred}{before the Maccabean period, around 165 B.C.; no individual psalm can be dated securely}{Before c. 165 B.C.}}%
}
% second-reading: Second Reading
\expandafter\def\csname triptychchronologyannotation@second-reading\endcsname{%
  \chronologyannotationgroup{composition}{preferred}{\textbf{Composition}: \chronologyannotationreach{1John.3.1}{1John}{true}\chronologyannotationreach{1John.3.2}{1John}{true}\chronologyannotationreach{1John.3.3}{1John}{true}\chronologyannotationclaim{composition.first-epistle-of-st-john}{composition}{catholic-traditional-v1}{preferred}{from the year 90 to 100 (approximately)}{c. A.D. 90--100}.}%
}
% gospel-acclamation: Gospel acclamation
\expandafter\def\csname triptychchronologyannotation@gospel-acclamation\endcsname{%
  \chronologyannotationgroup{composition}{disputed}{\textbf{Composition} -- disputed: \chronologyannotationreach{Matt.11.28}{Matt}{true}\chronologyannotationclaim{composition.gospel-of-matthew}{composition}{catholic-traditional-v1}{disputed}{about A.D. 38-45}{c. A.D. 38--45}; \chronologyannotationreach{Matt.11.28}{Matt}{true}\chronologyannotationclaim{composition.gospel-of-matthew}{composition}{catholic-traditional-v1}{disputed}{about the year 40-42}{c. A.D. 40--42}; \chronologyannotationreach{Matt.11.28}{Matt}{true}\chronologyannotationclaim{composition.gospel-of-matthew}{composition}{catholic-traditional-v1}{disputed}{the years 40-45}{A.D. 40--45}; \chronologyannotationreach{Matt.11.28}{Matt}{true}\chronologyannotationclaim{composition.gospel-of-matthew}{composition}{catholic-traditional-v1}{disputed}{about the year 60-68}{c. A.D. 60--68}; \chronologyannotationreach{Matt.11.28}{Matt}{true}\chronologyannotationclaim{composition.gospel-of-matthew}{composition}{catholic-traditional-v1}{disputed}{about the years 64-67}{c. A.D. 64--67}; \chronologyannotationreach{Matt.11.28}{Matt}{true}\chronologyannotationclaim{composition.gospel-of-matthew}{composition}{catholic-traditional-v1}{disputed}{about the year 50}{c. A.D. 50}.}%
}
% gospel: Gospel
\expandafter\def\csname triptychchronologyannotation@gospel\endcsname{%
  \chronologyannotationgroup{narrated-event}{research-pending}{\textbf{Event} -- No narrated-event date in the chronology corpus.}%
  \space \chronologyannotationgroup{composition}{disputed}{\textbf{Composition} -- disputed: \chronologyannotationreach{Matt.5.1}{Matt}{true}\chronologyannotationreach{Matt.5.10}{Matt}{true}\chronologyannotationreach{Matt.5.11}{Matt}{true}\chronologyannotationreach{Matt.5.12}{Matt}{true}\chronologyannotationreach{Matt.5.2}{Matt}{true}\chronologyannotationreach{Matt.5.3}{Matt}{true}\chronologyannotationreach{Matt.5.4}{Matt}{true}\chronologyannotationreach{Matt.5.5}{Matt}{true}\chronologyannotationreach{Matt.5.6}{Matt}{true}\chronologyannotationreach{Matt.5.7}{Matt}{true}\chronologyannotationreach{Matt.5.8}{Matt}{true}\chronologyannotationreach{Matt.5.9}{Matt}{true}\chronologyannotationclaim{composition.gospel-of-matthew}{composition}{catholic-traditional-v1}{disputed}{about A.D. 38-45}{c. A.D. 38--45}; \chronologyannotationreach{Matt.5.1}{Matt}{true}\chronologyannotationreach{Matt.5.10}{Matt}{true}\chronologyannotationreach{Matt.5.11}{Matt}{true}\chronologyannotationreach{Matt.5.12}{Matt}{true}\chronologyannotationreach{Matt.5.2}{Matt}{true}\chronologyannotationreach{Matt.5.3}{Matt}{true}\chronologyannotationreach{Matt.5.4}{Matt}{true}\chronologyannotationreach{Matt.5.5}{Matt}{true}\chronologyannotationreach{Matt.5.6}{Matt}{true}\chronologyannotationreach{Matt.5.7}{Matt}{true}\chronologyannotationreach{Matt.5.8}{Matt}{true}\chronologyannotationreach{Matt.5.9}{Matt}{true}\chronologyannotationclaim{composition.gospel-of-matthew}{composition}{catholic-traditional-v1}{disputed}{about the year 40-42}{c. A.D. 40--42}; \chronologyannotationreach{Matt.5.1}{Matt}{true}\chronologyannotationreach{Matt.5.10}{Matt}{true}\chronologyannotationreach{Matt.5.11}{Matt}{true}\chronologyannotationreach{Matt.5.12}{Matt}{true}\chronologyannotationreach{Matt.5.2}{Matt}{true}\chronologyannotationreach{Matt.5.3}{Matt}{true}\chronologyannotationreach{Matt.5.4}{Matt}{true}\chronologyannotationreach{Matt.5.5}{Matt}{true}\chronologyannotationreach{Matt.5.6}{Matt}{true}\chronologyannotationreach{Matt.5.7}{Matt}{true}\chronologyannotationreach{Matt.5.8}{Matt}{true}\chronologyannotationreach{Matt.5.9}{Matt}{true}\chronologyannotationclaim{composition.gospel-of-matthew}{composition}{catholic-traditional-v1}{disputed}{the years 40-45}{A.D. 40--45}; \chronologyannotationreach{Matt.5.1}{Matt}{true}\chronologyannotationreach{Matt.5.10}{Matt}{true}\chronologyannotationreach{Matt.5.11}{Matt}{true}\chronologyannotationreach{Matt.5.12}{Matt}{true}\chronologyannotationreach{Matt.5.2}{Matt}{true}\chronologyannotationreach{Matt.5.3}{Matt}{true}\chronologyannotationreach{Matt.5.4}{Matt}{true}\chronologyannotationreach{Matt.5.5}{Matt}{true}\chronologyannotationreach{Matt.5.6}{Matt}{true}\chronologyannotationreach{Matt.5.7}{Matt}{true}\chronologyannotationreach{Matt.5.8}{Matt}{true}\chronologyannotationreach{Matt.5.9}{Matt}{true}\chronologyannotationclaim{composition.gospel-of-matthew}{composition}{catholic-traditional-v1}{disputed}{about the year 60-68}{c. A.D. 60--68}; \chronologyannotationreach{Matt.5.1}{Matt}{true}\chronologyannotationreach{Matt.5.10}{Matt}{true}\chronologyannotationreach{Matt.5.11}{Matt}{true}\chronologyannotationreach{Matt.5.12}{Matt}{true}\chronologyannotationreach{Matt.5.2}{Matt}{true}\chronologyannotationreach{Matt.5.3}{Matt}{true}\chronologyannotationreach{Matt.5.4}{Matt}{true}\chronologyannotationreach{Matt.5.5}{Matt}{true}\chronologyannotationreach{Matt.5.6}{Matt}{true}\chronologyannotationreach{Matt.5.7}{Matt}{true}\chronologyannotationreach{Matt.5.8}{Matt}{true}\chronologyannotationreach{Matt.5.9}{Matt}{true}\chronologyannotationclaim{composition.gospel-of-matthew}{composition}{catholic-traditional-v1}{disputed}{about the years 64-67}{c. A.D. 64--67}; \chronologyannotationreach{Matt.5.1}{Matt}{true}\chronologyannotationreach{Matt.5.10}{Matt}{true}\chronologyannotationreach{Matt.5.11}{Matt}{true}\chronologyannotationreach{Matt.5.12}{Matt}{true}\chronologyannotationreach{Matt.5.2}{Matt}{true}\chronologyannotationreach{Matt.5.3}{Matt}{true}\chronologyannotationreach{Matt.5.4}{Matt}{true}\chronologyannotationreach{Matt.5.5}{Matt}{true}\chronologyannotationreach{Matt.5.6}{Matt}{true}\chronologyannotationreach{Matt.5.7}{Matt}{true}\chronologyannotationreach{Matt.5.8}{Matt}{true}\chronologyannotationreach{Matt.5.9}{Matt}{true}\chronologyannotationclaim{composition.gospel-of-matthew}{composition}{catholic-traditional-v1}{disputed}{about the year 50}{c. A.D. 50}.}%
}
% communion-antiphon: Communion antiphon
\expandafter\def\csname triptychchronologyannotation@communion-antiphon\endcsname{%
  \chronologyannotationgroup{composition}{disputed}{\textbf{Composition} -- disputed: \chronologyannotationreach{Matt.5.10}{Matt}{true}\chronologyannotationreach{Matt.5.8}{Matt}{true}\chronologyannotationreach{Matt.5.9}{Matt}{true}\chronologyannotationclaim{composition.gospel-of-matthew}{composition}{catholic-traditional-v1}{disputed}{about A.D. 38-45}{c. A.D. 38--45}; \chronologyannotationreach{Matt.5.10}{Matt}{true}\chronologyannotationreach{Matt.5.8}{Matt}{true}\chronologyannotationreach{Matt.5.9}{Matt}{true}\chronologyannotationclaim{composition.gospel-of-matthew}{composition}{catholic-traditional-v1}{disputed}{about the year 40-42}{c. A.D. 40--42}; \chronologyannotationreach{Matt.5.10}{Matt}{true}\chronologyannotationreach{Matt.5.8}{Matt}{true}\chronologyannotationreach{Matt.5.9}{Matt}{true}\chronologyannotationclaim{composition.gospel-of-matthew}{composition}{catholic-traditional-v1}{disputed}{the years 40-45}{A.D. 40--45}; \chronologyannotationreach{Matt.5.10}{Matt}{true}\chronologyannotationreach{Matt.5.8}{Matt}{true}\chronologyannotationreach{Matt.5.9}{Matt}{true}\chronologyannotationclaim{composition.gospel-of-matthew}{composition}{catholic-traditional-v1}{disputed}{about the year 60-68}{c. A.D. 60--68}; \chronologyannotationreach{Matt.5.10}{Matt}{true}\chronologyannotationreach{Matt.5.8}{Matt}{true}\chronologyannotationreach{Matt.5.9}{Matt}{true}\chronologyannotationclaim{composition.gospel-of-matthew}{composition}{catholic-traditional-v1}{disputed}{about the years 64-67}{c. A.D. 64--67}; \chronologyannotationreach{Matt.5.10}{Matt}{true}\chronologyannotationreach{Matt.5.8}{Matt}{true}\chronologyannotationreach{Matt.5.9}{Matt}{true}\chronologyannotationclaim{composition.gospel-of-matthew}{composition}{catholic-traditional-v1}{disputed}{about the year 50}{c. A.D. 50}.}%
}
% presentation: orientation-v3
\newcommand{\chronologyorientationrow}[2]{%
  \ifcsname triptychchronologyorientation@#1\endcsname
    \csname triptychchronologyorientation@#1\endcsname{#2}%
  \else
    \PackageError{triptych}{No chronology orientation row for #1}{}%
  \fi}
\expandafter\def\csname triptychchronologyorientation@row-df123084b173f7a9\endcsname#1{%
Responsorial Psalm: Ps. 24:1bc--2,3--4ab,5--6 (Heb.; Vg. 23), Cf. Ps. 24:6 (Heb.; Vg. 23) & Historical setting & Occasion date unresolved\\*[.1em]%
}
\expandafter\def\csname triptychchronologyorientation@row-64506719eac72d70\endcsname#1{%
 & Composition: Psalms & \chronologyannotationreach{Ps.23.1}{Ps}{true}\chronologyannotationreach{Ps.23.2}{Ps}{true}\chronologyannotationreach{Ps.23.3}{Ps}{true}\chronologyannotationreach{Ps.23.4}{Ps}{true}\chronologyannotationreach{Ps.23.5}{Ps}{true}\chronologyannotationreach{Ps.23.6}{Ps}{true}\chronologyannotationclaim{critical.psalms.latest-composition-boundary}{composition}{catholic-critical-v1}{preferred}{before the Maccabean period, around 165 B.C.; no individual psalm can be dated securely}{Before c.~165 B.C.}\\[.3em]%
}
\expandafter\def\csname triptychchronologyorientation@row-f88fad136044987c\endcsname#1{%
Gospel: Matt. 5:1--12a & Narrated event: Jesus teaches the Beatitudes on an unnamed mountain in Galilee & Date unresolved\\*[.1em]%
}
\expandafter\def\csname triptychchronologyorientation@row-8eebef4dea81e70b\endcsname#1{%
 & Composition: Matthew; Semitic original associated with Palestine & Disputed: \chronologyannotationreach{Matt.5.1}{Matt}{true}\chronologyannotationreach{Matt.5.10}{Matt}{true}\chronologyannotationreach{Matt.5.11}{Matt}{true}\chronologyannotationreach{Matt.5.12}{Matt}{true}\chronologyannotationreach{Matt.5.2}{Matt}{true}\chronologyannotationreach{Matt.5.3}{Matt}{true}\chronologyannotationreach{Matt.5.4}{Matt}{true}\chronologyannotationreach{Matt.5.5}{Matt}{true}\chronologyannotationreach{Matt.5.6}{Matt}{true}\chronologyannotationreach{Matt.5.7}{Matt}{true}\chronologyannotationreach{Matt.5.8}{Matt}{true}\chronologyannotationreach{Matt.5.9}{Matt}{true}\chronologyannotationclaim{composition.gospel-of-matthew}{composition}{catholic-traditional-v1}{disputed}{about A.D. 38-45}{c.~A.D.~38--45}; \chronologyannotationreach{Matt.5.1}{Matt}{true}\chronologyannotationreach{Matt.5.10}{Matt}{true}\chronologyannotationreach{Matt.5.11}{Matt}{true}\chronologyannotationreach{Matt.5.12}{Matt}{true}\chronologyannotationreach{Matt.5.2}{Matt}{true}\chronologyannotationreach{Matt.5.3}{Matt}{true}\chronologyannotationreach{Matt.5.4}{Matt}{true}\chronologyannotationreach{Matt.5.5}{Matt}{true}\chronologyannotationreach{Matt.5.6}{Matt}{true}\chronologyannotationreach{Matt.5.7}{Matt}{true}\chronologyannotationreach{Matt.5.8}{Matt}{true}\chronologyannotationreach{Matt.5.9}{Matt}{true}\chronologyannotationclaim{composition.gospel-of-matthew}{composition}{catholic-traditional-v1}{disputed}{about the year 40-42}{c.~A.D.~40--42}; \chronologyannotationreach{Matt.5.1}{Matt}{true}\chronologyannotationreach{Matt.5.10}{Matt}{true}\chronologyannotationreach{Matt.5.11}{Matt}{true}\chronologyannotationreach{Matt.5.12}{Matt}{true}\chronologyannotationreach{Matt.5.2}{Matt}{true}\chronologyannotationreach{Matt.5.3}{Matt}{true}\chronologyannotationreach{Matt.5.4}{Matt}{true}\chronologyannotationreach{Matt.5.5}{Matt}{true}\chronologyannotationreach{Matt.5.6}{Matt}{true}\chronologyannotationreach{Matt.5.7}{Matt}{true}\chronologyannotationreach{Matt.5.8}{Matt}{true}\chronologyannotationreach{Matt.5.9}{Matt}{true}\chronologyannotationclaim{composition.gospel-of-matthew}{composition}{catholic-traditional-v1}{disputed}{the years 40-45}{A.D.~40--45}; \chronologyannotationreach{Matt.5.1}{Matt}{true}\chronologyannotationreach{Matt.5.10}{Matt}{true}\chronologyannotationreach{Matt.5.11}{Matt}{true}\chronologyannotationreach{Matt.5.12}{Matt}{true}\chronologyannotationreach{Matt.5.2}{Matt}{true}\chronologyannotationreach{Matt.5.3}{Matt}{true}\chronologyannotationreach{Matt.5.4}{Matt}{true}\chronologyannotationreach{Matt.5.5}{Matt}{true}\chronologyannotationreach{Matt.5.6}{Matt}{true}\chronologyannotationreach{Matt.5.7}{Matt}{true}\chronologyannotationreach{Matt.5.8}{Matt}{true}\chronologyannotationreach{Matt.5.9}{Matt}{true}\chronologyannotationclaim{composition.gospel-of-matthew}{composition}{catholic-traditional-v1}{disputed}{about the year 60-68}{c.~A.D.~60--68}; \chronologyannotationreach{Matt.5.1}{Matt}{true}\chronologyannotationreach{Matt.5.10}{Matt}{true}\chronologyannotationreach{Matt.5.11}{Matt}{true}\chronologyannotationreach{Matt.5.12}{Matt}{true}\chronologyannotationreach{Matt.5.2}{Matt}{true}\chronologyannotationreach{Matt.5.3}{Matt}{true}\chronologyannotationreach{Matt.5.4}{Matt}{true}\chronologyannotationreach{Matt.5.5}{Matt}{true}\chronologyannotationreach{Matt.5.6}{Matt}{true}\chronologyannotationreach{Matt.5.7}{Matt}{true}\chronologyannotationreach{Matt.5.8}{Matt}{true}\chronologyannotationreach{Matt.5.9}{Matt}{true}\chronologyannotationclaim{composition.gospel-of-matthew}{composition}{catholic-traditional-v1}{disputed}{about the years 64-67}{c.~A.D.~64--67}; \chronologyannotationreach{Matt.5.1}{Matt}{true}\chronologyannotationreach{Matt.5.10}{Matt}{true}\chronologyannotationreach{Matt.5.11}{Matt}{true}\chronologyannotationreach{Matt.5.12}{Matt}{true}\chronologyannotationreach{Matt.5.2}{Matt}{true}\chronologyannotationreach{Matt.5.3}{Matt}{true}\chronologyannotationreach{Matt.5.4}{Matt}{true}\chronologyannotationreach{Matt.5.5}{Matt}{true}\chronologyannotationreach{Matt.5.6}{Matt}{true}\chronologyannotationreach{Matt.5.7}{Matt}{true}\chronologyannotationreach{Matt.5.8}{Matt}{true}\chronologyannotationreach{Matt.5.9}{Matt}{true}\chronologyannotationclaim{composition.gospel-of-matthew}{composition}{catholic-traditional-v1}{disputed}{about the year 50}{c.~A.D.~50}\\[.3em]%
}
\expandafter\def\csname triptychchronologyorientation@row-e33616248ec90765\endcsname#1{%
Communion antiphon: Matt. 5:8--10 & Composition: Matthew; Semitic original associated with Palestine & Disputed: \chronologyannotationreach{Matt.5.10}{Matt}{true}\chronologyannotationreach{Matt.5.8}{Matt}{true}\chronologyannotationreach{Matt.5.9}{Matt}{true}\chronologyannotationclaim{composition.gospel-of-matthew}{composition}{catholic-traditional-v1}{disputed}{about A.D. 38-45}{c.~A.D.~38--45}; \chronologyannotationreach{Matt.5.10}{Matt}{true}\chronologyannotationreach{Matt.5.8}{Matt}{true}\chronologyannotationreach{Matt.5.9}{Matt}{true}\chronologyannotationclaim{composition.gospel-of-matthew}{composition}{catholic-traditional-v1}{disputed}{about the year 40-42}{c.~A.D.~40--42}; \chronologyannotationreach{Matt.5.10}{Matt}{true}\chronologyannotationreach{Matt.5.8}{Matt}{true}\chronologyannotationreach{Matt.5.9}{Matt}{true}\chronologyannotationclaim{composition.gospel-of-matthew}{composition}{catholic-traditional-v1}{disputed}{the years 40-45}{A.D.~40--45}; \chronologyannotationreach{Matt.5.10}{Matt}{true}\chronologyannotationreach{Matt.5.8}{Matt}{true}\chronologyannotationreach{Matt.5.9}{Matt}{true}\chronologyannotationclaim{composition.gospel-of-matthew}{composition}{catholic-traditional-v1}{disputed}{about the year 60-68}{c.~A.D.~60--68}; \chronologyannotationreach{Matt.5.10}{Matt}{true}\chronologyannotationreach{Matt.5.8}{Matt}{true}\chronologyannotationreach{Matt.5.9}{Matt}{true}\chronologyannotationclaim{composition.gospel-of-matthew}{composition}{catholic-traditional-v1}{disputed}{about the years 64-67}{c.~A.D.~64--67}; \chronologyannotationreach{Matt.5.10}{Matt}{true}\chronologyannotationreach{Matt.5.8}{Matt}{true}\chronologyannotationreach{Matt.5.9}{Matt}{true}\chronologyannotationclaim{composition.gospel-of-matthew}{composition}{catholic-traditional-v1}{disputed}{about the year 50}{c.~A.D.~50}\\[.3em]%
}
\expandafter\def\csname triptychchronologyorientation@row-d4d8a04526f92512\endcsname#1{%
Gospel acclamation: Matt. 11:28 & Composition: Matthew; Semitic original associated with Palestine & Disputed: \chronologyannotationreach{Matt.11.28}{Matt}{true}\chronologyannotationclaim{composition.gospel-of-matthew}{composition}{catholic-traditional-v1}{disputed}{about A.D. 38-45}{c.~A.D.~38--45}; \chronologyannotationreach{Matt.11.28}{Matt}{true}\chronologyannotationclaim{composition.gospel-of-matthew}{composition}{catholic-traditional-v1}{disputed}{about the year 40-42}{c.~A.D.~40--42}; \chronologyannotationreach{Matt.11.28}{Matt}{true}\chronologyannotationclaim{composition.gospel-of-matthew}{composition}{catholic-traditional-v1}{disputed}{the years 40-45}{A.D.~40--45}; \chronologyannotationreach{Matt.11.28}{Matt}{true}\chronologyannotationclaim{composition.gospel-of-matthew}{composition}{catholic-traditional-v1}{disputed}{about the year 60-68}{c.~A.D.~60--68}; \chronologyannotationreach{Matt.11.28}{Matt}{true}\chronologyannotationclaim{composition.gospel-of-matthew}{composition}{catholic-traditional-v1}{disputed}{about the years 64-67}{c.~A.D.~64--67}; \chronologyannotationreach{Matt.11.28}{Matt}{true}\chronologyannotationclaim{composition.gospel-of-matthew}{composition}{catholic-traditional-v1}{disputed}{about the year 50}{c.~A.D.~50}\\[.3em]%
}
\expandafter\def\csname triptychchronologyorientation@row-f480225fc28156b1\endcsname#1{%
Second Reading: 1 John. 3:1--3 & Composition: First John; Ephesus probable & \chronologyannotationreach{1John.3.1}{1John}{true}\chronologyannotationreach{1John.3.2}{1John}{true}\chronologyannotationreach{1John.3.3}{1John}{true}\chronologyannotationclaim{composition.first-epistle-of-st-john}{composition}{catholic-traditional-v1}{preferred}{from the year 90 to 100 (approximately)}{c.~A.D.~90--100}\\[.3em]%
}
\expandafter\def\csname triptychchronologyorientation@row-38dd8f1a45d0b6f6\endcsname#1{%
First Reading: Apoc. 7:2--4,9--14 & Composition: Apocalypse; seer on Patmos; heavenly scenes are visions & Disputed: \chronologyannotationreach{Apoc.7.10}{Apoc}{true}\chronologyannotationreach{Apoc.7.11}{Apoc}{true}\chronologyannotationreach{Apoc.7.12}{Apoc}{true}\chronologyannotationreach{Apoc.7.13}{Apoc}{true}\chronologyannotationreach{Apoc.7.14}{Apoc}{true}\chronologyannotationreach{Apoc.7.2}{Apoc}{true}\chronologyannotationreach{Apoc.7.3}{Apoc}{true}\chronologyannotationreach{Apoc.7.4}{Apoc}{true}\chronologyannotationreach{Apoc.7.9}{Apoc}{true}\chronologyannotationclaim{composition.apocalypse}{composition}{catholic-traditional-v1}{disputed}{the reign of Domitian ... A.D. 81-96}{A.D.~81--96}; \chronologyannotationreach{Apoc.7.10}{Apoc}{true}\chronologyannotationreach{Apoc.7.11}{Apoc}{true}\chronologyannotationreach{Apoc.7.12}{Apoc}{true}\chronologyannotationreach{Apoc.7.13}{Apoc}{true}\chronologyannotationreach{Apoc.7.14}{Apoc}{true}\chronologyannotationreach{Apoc.7.2}{Apoc}{true}\chronologyannotationreach{Apoc.7.3}{Apoc}{true}\chronologyannotationreach{Apoc.7.4}{Apoc}{true}\chronologyannotationreach{Apoc.7.9}{Apoc}{true}\chronologyannotationclaim{composition.apocalypse}{composition}{catholic-traditional-v1}{disputed}{from the year 90 to 100 (approximately)}{c.~A.D.~90--100}\\[.3em]%
}

% Generated. Do not edit.
% Complete orientation-v3 table; wording, dates and order belong to the corpus.
\begin{chronologysharedtable}
\chronologyorientationrow{row-df123084b173f7a9}{}
\chronologyorientationrow{row-64506719eac72d70}{}
\midrule
\chronologyorientationrow{row-f88fad136044987c}{}
\chronologyorientationrow{row-8eebef4dea81e70b}{}
\midrule
\chronologyorientationrow{row-e33616248ec90765}{}
\midrule
\chronologyorientationrow{row-d4d8a04526f92512}{}
\midrule
\chronologyorientationrow{row-f480225fc28156b1}{}
\midrule
\chronologyorientationrow{row-38dd8f1a45d0b6f6}{}
\end{chronologysharedtable}
{\footnotesize Psalm numbering: Vg. = Vulgate; Heb. = Hebrew.\par}


\subsection{Evidence for each date}
\label{sec:date-location-evidence}
{\small For each passage, as the Scriptural Date and Location table names it: what is dated, how the corpus weighs each answer, the source's own words, and the sources cited. A subject the corpus does not date is reported as such and given no date.\par}
{\small
\par\smallskip\noindent\textit{Responsorial Psalm: Ps. 24:1bc--2,3--4ab,5--6 (Heb.; Vg. 23), Cf. Ps. 24:6 (Heb.; Vg. 23)}\par\nobreak
\begin{itemize}[leftmargin=1.2em,itemsep=1pt,topsep=1pt,parsep=0pt]
\item \textbf{Historical setting}. No date: the chronology corpus has not yet dated this (research pending).
\item \textbf{Composition} (preferred). Psalms: Before c.~165 B.C. (``before the Maccabean period, around 165 B.C.; no individual psalm can be dated securely''; Catholic critical chronology; New American Bible, Revised Edition --- USCCB web state (2026), Psalms introduction).
\end{itemize}
\par\smallskip\noindent\textit{Gospel: Matt. 5:1--12a}\par\nobreak
\begin{itemize}[leftmargin=1.2em,itemsep=1pt,topsep=1pt,parsep=0pt]
\item \textbf{Narrated event}. Jesus teaches the Beatitudes on an unnamed mountain in Galilee (The Holy Bible, Douay-Rheims (Challoner revision), Project Gutenberg eBook 1581: Matt. 4:23, 4:25, 5:1, 5:2, 5:3, 5:12). No date: the chronology corpus has not yet dated this (research pending).
\item \textbf{Composition} (disputed). The Gospel of St Matthew: c.~A.D.~38--45 (``about A.D. 38-45''; The Catholic Encyclopedia, Volume 10 (1911); The Catholic Encyclopedia, Volume 14 (New York, 1912)); c.~A.D.~40--42 (``about the year 40-42''; The Catholic Encyclopedia, Volume 10 (1911)); A.D.~40--45 (``the years 40-45''; The Catholic Encyclopedia, Volume 10 (1911)); c.~A.D.~60--68 (``about the year 60-68''; The Catholic Encyclopedia, Volume 10 (1911)); c.~A.D.~64--67 (``about the years 64-67''; The Catholic Encyclopedia, Volume 10 (1911)); c.~A.D.~50 (``about the year 50''; The Catholic Encyclopedia, Volume 14 (New York, 1912)). Place: Semitic original associated with Palestine (The Catholic Encyclopedia, Volume 10 (1911)).
\end{itemize}
\par\smallskip\noindent\textit{Communion antiphon: Matt. 5:8--10}\par\nobreak
\begin{itemize}[leftmargin=1.2em,itemsep=1pt,topsep=1pt,parsep=0pt]
\item \textbf{Composition}. As under Gospel: Matt. 5:1--12a.
\end{itemize}
\par\smallskip\noindent\textit{Gospel acclamation: Matt. 11:28}\par\nobreak
\begin{itemize}[leftmargin=1.2em,itemsep=1pt,topsep=1pt,parsep=0pt]
\item \textbf{Composition}. As under Gospel: Matt. 5:1--12a.
\end{itemize}
\par\smallskip\noindent\textit{Second Reading: 1 John. 3:1--3}\par\nobreak
\begin{itemize}[leftmargin=1.2em,itemsep=1pt,topsep=1pt,parsep=0pt]
\item \textbf{Composition} (preferred). The First Epistle of St John: c.~A.D.~90--100 (``from the year 90 to 100 (approximately)''; The Catholic Encyclopedia, Volume 14 (New York, 1912); The Catholic Encyclopedia, Volume 8 (New York, 1910)). Place: Ephesus probable (The Catholic Encyclopedia, Volume 8 (New York, 1910)).
\end{itemize}
\par\smallskip\noindent\textit{First Reading: Apoc. 7:2--4,9--14}\par\nobreak
\begin{itemize}[leftmargin=1.2em,itemsep=1pt,topsep=1pt,parsep=0pt]
\item \textbf{Composition} (disputed). The Apocalypse of St John: A.D.~81--96 (``the reign of Domitian ... A.D. 81-96''; The Catholic Encyclopedia, Volume 1 (New York, 1907)); c.~A.D.~90--100 (``from the year 90 to 100 (approximately)''; The Catholic Encyclopedia, Volume 14 (New York, 1912)). Place: seer on Patmos; heavenly scenes are visions (The Holy Bible, Douay-Rheims (Challoner revision), Project Gutenberg eBook 1581: Apoc. 1:9, 4:1, 7:9).
\end{itemize}
\par}
